\documentclass[11pt]{article}
\usepackage[margin=0.75in]{geometry}
\usepackage{graphicx}
\usepackage{float}
\usepackage{booktabs}
\usepackage{rotating}
\setlength{\parindent}{0pt}
\setlength{\parskip}{0.5em}
\begin{document}
\title{GEOS-MPM Examples Suite Report}
\author{Generated by examples/runExampleSuite}
\date{\today}
\maketitle
\section*{Running and rerunning}
Each example directory contains a copyable \texttt{pfw\_input\_*.py} and a \texttt{runProblem} wrapper.  The wrapper stages an isolated case directory under \texttt{testRunDirectory}, copies the needed PFW files, runs the particle-file writer, submits GEOS, submits dependent post-processing, and copies figures/logs back to the example \texttt{output/} directory.
\begin{verbatim}
cd scripts/preProcessing/materialPointMethodParticleFileWriter/examples
./runExampleSuite all --python /usr/tce/bin/python3 --force
./runExampleSuite all --case borehole_Ghareb --python /usr/tce/bin/python3 --force
./runExampleSuite report --python /usr/tce/bin/python3
\end{verbatim}
If output exists, \texttt{runProblem} prompts before overwriting. Use \texttt{--force} for noninteractive reruns.
\section*{Timing summary}
\begin{table}[H]
\centering
\scriptsize
\begin{tabular}{llrrrrrr}
\toprule
Example & Family & GEOS job & GEOS state & GEOS wall & GEOS CPU & Post wall & Post CPU \\
\midrule
brazilianDisk\_BSpline & Brazilian disk & 6044292 & COMPLETED & 00:02:18 & 00:00:00 & 00:00:10 & 00:00:00 \\
\bottomrule
\end{tabular}
\end{table}
\clearpage
\section*{brazilianDisk\_BSpline}
Brazilian disk compression using the new SinglePointBSpline particle type.  The case keeps the FMPM update path and quartz damage material but uses particleType=1 so GEOS exercises ParticleType::SinglePointBSpline.
\begin{figure}[H]
\centering
\setlength{\unitlength}{0.0045\textwidth}\begin{picture}(120,82)\put(12,12){\framebox(96,58){}}\put(60,41){\circle{38}}\put(44,41){\vector(0,-1){18}}\put(76,41){\vector(0,1){18}}\put(25,73){moving y+}\put(25,3){moving y-}\put(6,31){\rotatebox{90}{periodic x}}\put(109,31){\rotatebox{-90}{periodic x}}\end{picture}
\caption{Schematic for brazilianDisk\_BSpline.}
\end{figure}
\textbf{Code features exercised.}
\begin{itemize}
\item SinglePointBSpline particle type
\item FMPM update
\item quartz damage material
\item moving y boundaries
\item periodic x boundaries
\item reaction history
\end{itemize}
\begin{figure}[H]
\centering
\includegraphics[width=0.31\textwidth]{../brazilianDisk_BSpline/output/brazilianDisk_BSpline_initial_particleStrengthScale.png}
\includegraphics[width=0.31\textwidth]{../brazilianDisk_BSpline/output/brazilianDisk_BSpline_final_particleDamage.png}
\includegraphics[width=0.31\textwidth]{../brazilianDisk_BSpline/output/brazilianDisk_BSpline_reactions_y.png}
\caption{Initial state, final state, and reaction history for brazilianDisk\_BSpline.}
\end{figure}
\end{document}
